\documentclass[11pt]{article}
% common.tex — minimal noise setup
% ----------------------------

% Silence package messages before anything else
\RequirePackage{silence}

% Completely hide "info" messages from LaTeX packages
\WarningFilter*{latex}{}
\WarningFilter*{hyperref}{}
\WarningFilter*{pgf}{}
\WarningFilter*{tikz}{}

% Optionally: hide some hyperref PDF string warnings that are common but harmless
\WarningFilter{hyperref}{Token not allowed in a PDF string}

% ----------------------------
% Now load the usual packages
\usepackage[utf8]{inputenc}
\usepackage[T1]{fontenc}
\usepackage{amsmath, amssymb, amsthm, amsfonts}
\usepackage{geometry}
\usepackage{hyperref}
\usepackage{graphicx}
\usepackage{physics}
\usepackage{booktabs}
\usepackage{listings}
\usepackage{bookmark}
\usepackage{amscd}
\usepackage{tikz-cd}
\usepackage{mathtools}
\usepackage{xspace}
\usepackage[backend=biber, style=numeric]{biblatex}
\usepackage{glossaries}
\usepackage{tcolorbox}

% ----------------------------
% Theorems, macros
\newtheorem{lemma}{Lemma}
\newcommand{\IaMe}{\(IaM^e\)\xspace}
\newcommand{\pdflink}[1]{}

% ----------------------------
% Page layout
\geometry{letterpaper, margin=1in}
\setlength{\parindent}{0pt}
\setlength{\parskip}{1em}

% Hyperref settings
\hypersetup{
    colorlinks=true,
    linkcolor=black,
    citecolor=black,
    filecolor=magenta,
    urlcolor=blue
}

% TOC formatting
\usepackage{tocloft}
\setlength{\cftbeforechapskip}{1.0em}
\setlength{\cftbeforesecskip}{0pt}
\setlength{\cftbeforesubsecskip}{0pt}
\renewcommand{\cftchapleader}{\cftdotfill{\cftdotsep}}

% Glossary
\defglsentryfmt{\ifglsused{\glslabel}{\glsgenentryfmt}{\textit{\glsgenentryfmt}}}

% Author
\author{$IaM^e$}


\input{../glossary.tex}
\addbibresource{../references.bib}

\title{Why Something Rather Than Nothing}
  
\date{2008}

\begin{document}

\maketitle

\begin{abstract}
Raw informational structures possess no intrinsic semantics.
Any assignment of meaning, ontology, mathematics, or physical
interpretation arises only relative to an interpretive framework.

Consequently, ontological claims that depend purely on representational
choice are gauge-dependent and cannot possess observer-independent
metaphysical significance.

The following formal framework develops the consequences of this principle
for the classical question ``Why something rather than nothing?''
\end{abstract}


\section{Introduction}

The question of why anything exists at all has long been considered a fundamental mystery of metaphysics. By treating existence and non-existence through the lens of informational structures and gauge invariance, we can analyze the structural properties of this problem mathematically. 

In this paper, we formalize this approach by treating representational choices as symmetries. Section~\ref{sec:igit} outlines the core axioms of this Interpretive Gauge Invariance Thesis and demonstrates how the classical question ultimately dissolves under observer self-location.

\section{Interpretive Gauge Invariance Thesis}
\label{sec:igit}

\begin{definition}
A \emph{binary distinction} is a pair of mutually exclusive and jointly exhaustive states, canonically represented by the set $\mathbf{2} = \{0, 1\}$.
\end{definition}

\begin{axiom}[Binary Representability of Existence]
The distinction between existence and non-existence (``to be or not to be'') can be represented as a binary distinction: one state corresponding to the empty set $\varnothing$ (``nothing'') and the other to a non-empty structure (``something'').
\end{axiom}

\begin{axiom}[Semantic Non-Intrinsicality]
No raw informational structure possesses an intrinsic interpretation.
Any semantics assigned to a representation arises only relative
to an interpretation.
\end{axiom}

\begin{axiom}[Closure under Set-Theoretic Construction]
Arbitrary set-theoretic objects (power sets, function spaces $\mathbf{2}^S$, relations, etc.) may be constructed from binary distinctions. All such constructions remain subject to Axioms 1 and 2.
\end{axiom}

\begin{axiom}[Plenitude]
Every internally consistent mathematical structure possesses equal ontological standing. There is no external ``existence predicate'' that selects one consistent structure as actual over any other.
\end{axiom}

\begin{theorem}[Representational Symmetry]
Under Axioms 1--2, the binary encoding of ``something'' versus ``nothing'' does not confer ontological priority on either side. The choice of which label designates the empty configuration is arbitrary.
\end{theorem}

\begin{theorem}[No Privileged Configuration via Labeling]
Let $S$ be any index set and let $\mathbf{2}^S$ be the space of all possible configurations (functions $f: S \to \{0,1\}$).  

No configuration $f \in \mathbf{2}^S$ --- including the all-zero configuration (the empty universe) and any information-rich configuration containing observers --- is distinguished from the others \emph{merely by representational labeling conventions}. 
\end{theorem}

\begin{proof}
Consider the boolean inversion operator $\neg: \mathbf{2} \to \mathbf{2}$. This induces a bijection $\Phi: \mathbf{2}^S \to \mathbf{2}^S$ by composition, $\Phi(f) = \neg \circ f$. By Axiom 2, $\Phi$ preserves all purely ontological standing. Since the all-zero configuration maps to the all-one configuration under $\Phi$, any predicate assigning ``actuality'' based solely on the label assignment is gauge-dependent and physically non-invariant.
\end{proof}

\begin{theorem}[Observer Self-Location]
Any self-aware observer can only exist within those configurations in $\mathbf{2}^S$ that contain the necessary informational structure to support consciousness and consistent memories. Observers therefore necessarily self-locate within observer-compatible configurations. there is no requirement to explain a causal or temporal transition from an empty configuration to a non-empty one.
\end{theorem}

\begin{corollary}[Dissolution of the Classical Question]
The question ``Why something rather than nothing?'' presupposes that the empty configuration possesses a special metaphysical status that requires explanation for why it was not realized.  

Under Axioms 1--4, this presupposition fails:  
\begin{itemize}
    \item The distinction is representable as a binary choice whose labels are arbitrary (Axioms 1--2);
    \item All consistent configurations have equal ontological standing (Axiom 4);
    \item Observers can only appear in structured configurations (Theorem 3).
\end{itemize}
The classical question is therefore ill-posed. It is replaced by the milder issue of observer self-location within the ensemble of consistent mathematical structures.
\end{corollary}

\begin{remark}
This framework distinguishes three notions:
\begin{enumerate}
    \item \emph{Representational symmetry} (Axiom 2): labels are arbitrary;
    \item \emph{Structural properties}: different configurations generally have different internal structure (the empty configuration is not isomorphic to a rich one);
    \item \emph{Ontological standing} (Axiom 4): all internally consistent structures are on equal footing.
\end{enumerate}
The argument relies on the move from (1) and (3) to the dissolution, while fully acknowledging (2).
\end{remark}

This formalization shows that once one accepts the neutrality of representational labels and the equal ontological status of consistent structures, the traditional asymmetry between ``something'' and ``nothing'' disappears.

\end{document}
